\documentclass[11pt]{article}
\usepackage{amsmath,amssymb,amsfonts}
\usepackage{geometry}
\geometry{margin=2.5cm}

\title{Fonctions de plusieurs variables}
\author{}
\date{}

\begin{document}
\maketitle

\section{Introduction}

La plupart des relations en ingénierie et en physique comportent plusieurs paramètres.
On aimerait comprendre le comportement d’un système en fonction de ces paramètres.

\subsection*{Exemples}

\begin{enumerate}
\item \textbf{Boîte de conserve}

On considère une boîte cylindrique de rayon $r$ et de hauteur $h$.

La surface extérieure est :
\[
S = 2\pi r^2 + 2\pi r h.
\]

Les variables ou paramètres sont $r$ et $h$.

Quelles valeurs pour $r$ et $h$ afin que $S$ soit minimale, sous la contrainte :
\[
V = \pi r^2 h = \text{cste}.
\]

\textbf{Réponse :} $r = h$.

\item \textbf{Loi des gaz parfaits}

La loi des gaz parfaits est :
\[
PV = nRT.
\]

On peut écrire :
\[
P = \frac{nRT}{V} = f(n,T,V).
\]

\end{enumerate}

L’idée est de généraliser les notions connues pour les fonctions à une variable :
limites, continuité, dérivées, intégrales.

\section{Limites}

Soit la fonction :
\[
f(x,y) = \frac{xy}{x+y}.
\]

\subsection*{Domaine de définition}

Le domaine de définition est :
\[
D_f : x+y \neq 0 \quad \Leftrightarrow \quad y \neq -x.
\]

Le domaine de définition est un sous-ensemble de $\mathbb{R}^2$ défini par une équation.

\subsection*{Limite en $(0,0)$}

Quelle est la limite de $f$ lorsque $x$ tend vers $0$ et $y$ tend vers $0$, c’est-à-dire
lorsque $(x,y) \to (0,0)$ ?

La fonction $f(0,0)$ n’est pas définie.

On a une infinité de façons de tendre vers $0$.
Le point $M(x,y)$ se rapproche de l’origine $O$.

La limite dépend-elle de la façon dont on se rapproche de $0$ ?

On se place en coordonnées polaires :
\[
\begin{cases}
x = r\cos\theta, \\
y = r\sin\theta,
\end{cases}
\qquad r = \sqrt{x^2 + y^2}, \quad r \ge 0.
\]

Alors :
\[
f(x,y) = g(r,\theta)
= \frac{r\cos\theta \cdot r\sin\theta}{r\cos\theta + r\sin\theta}
= \frac{r\sin\theta\cos\theta}{\cos\theta + \sin\theta}.
\]

\[
(x,y) \to (0,0)
\quad \Longleftrightarrow \quad
\begin{cases}
r \to 0, \\
\forall \theta \in \mathbb{R}.
\end{cases}
\]

On a :
\[
\lim_{r \to 0} g(r,\theta) = 0
\quad \Longrightarrow \quad
\lim_{(x,y)\to(0,0)} f(x,y) = 0.
\]

\subsection*{Exemple où la limite n’existe pas}

Soit la fonction :
\[
h(x,y) = \frac{xy}{x^2 + y^2}.
\]

En coordonnées polaires :
\[
h(r,\theta)
= \frac{r^2\sin\theta\cos\theta}{r^2(\cos^2\theta + \sin^2\theta)}
= \sin\theta\cos\theta
= \frac{1}{2}\sin(2\theta).
\]

Ainsi :
\[
\lim_{r \to 0} h(r,\theta) = \frac{1}{2}\sin(2\theta),
\]
ce qui dépend de la direction d’approche.

La limite en $(0,0)$ dépend donc de la direction d’approche :
on a une infinité de limites, et la limite n’existe pas.

\section{Dérivées et différentielles}

Une \og fonction de plusieurs variables \fg{} peut être :

\begin{itemize}
\item une fonction scalaire sur un champ scalaire :
\[
f : \mathbb{R}^n \to \mathbb{R},
\quad (x_1,\ldots,x_n) \mapsto f(x_1,\ldots,x_n) \in \mathbb{R}.
\]
\end{itemize}

\textbf{Exemple :}
\[
f(x,y,z) = x^2yz,
\quad f : \mathbb{R}^3 \to \mathbb{R}.
\]

\end{document}

